\section*{Bab \ref{ch:g10:vectors} --- Vektor}

\begin{solution}{exo:g10:vectors:1}
$\vect{AB} + \vect{BD} = \vect{AD}$ (Chasles).

$\vect{KL} + \vect{LM} + \vect{MK} = \vect{KK} = \vec 0$.

$\vect{AC} - \vect{BC} = \vect{AC} + \vect{CB} = \vect{AB}$.

$\vect{AB} + \vect{CA} = \vect{CA} + \vect{AB} = \vect{CB}$.
\end{solution}

\begin{solution}{exo:g10:vectors:2}
$\vect{AB}\,(5-2,\ 3-1) = (3, 2)$; \quad
$\vect{BC}\,(-1-5,\ 4-3) = (-6, 1)$; \quad
$\vect{AC}\,(-1-2,\ 4-1) = (-3, 3)$.
Periksa: $(3, 2) + (-6, 1) = (-3, 3)$, \omterm{def:g10:coordgeom:system}{koordinat} demi \omterm{def:g10:coordgeom:system}{koordinat}.
\end{solution}

\begin{solution}{exo:g10:vectors:3}
$\vec u + \vec v = (2 + (-1),\ -3 + 4) = (1, 1)$.

$3\vec u = (6, -9)$.

$2\vec u - \vec v = (4 - (-1),\ -6 - 4) = (5, -10)$.

$\norm{\vec u} = \sqrt{2^2 + (-3)^2} = \sqrt{13}$.
\end{solution}

\begin{solution}{exo:g10:vectors:4}
(a) $4 \times 9 - 6 \times 6 = 36 - 36 = 0$: \omterm{def:g10:vectors:collinear}{kolinear}
($\vec v = \frac32 \vec u$).

(b) $2 \times 10 - (-5) \times (-4) = 20 - 20 = 0$: \omterm{def:g10:vectors:collinear}{kolinear}
($\vec v = -2\vec u$).

(c) $3 \times 3 - 1 \times 1 = 8 \neq 0$: tidak \omterm{def:g10:vectors:collinear}{kolinear}.
\end{solution}

\begin{solution}{exo:g10:vectors:5}
$\vect{AB}\,(4, -2)$, dan $\vect{DC}\,(6 - x_D,\ 3 - y_D)$. Syarat
$\vect{AB} = \vect{DC}$ memberi $6 - x_D = 4$ dan $3 - y_D = -2$, jadi
$D(2, 5)$.
\end{solution}

\begin{solution}{exo:g10:vectors:6}
\emph{1.} $\vect{AB}\,(3, 2)$ dan $\vect{AC}\,(6, -2)$, sehingga
$\vect{AB} + \vect{AC} = (9, 0)$ dan
$M = (-2 + 9,\ 1 + 0) = (7, 1)$.

\emph{2.} Menurut aturan jajargenjang, $\vect{AM} = \vect{AB} + \vect{AC}$
berarti $ABMC$ adalah jajargenjang ($M$ titik sudut keempat yang berhadapan
dengan $A$). Dapat diperiksa: $\vect{AB}\,(3,2)$ dan
$\vect{CM}\,(7-4,\ 1-(-1)) = (3, 2)$ memang sama.
\end{solution}

\begin{solution}{exo:g10:vectors:7}
Tripel pertama: $\vect{AB}\,(3, 2)$, $\vect{AC}\,(9, 6)$;
$3 \times 6 - 2 \times 9 = 0$: segaris ($\vect{AC} = 3\vect{AB}$).

Tripel kedua: $\vect{AB}\,(2, -1)$, $\vect{AC}\,(8, -4)$;
$2 \times (-4) - (-1) \times 8 = -8 + 8 = 0$: segaris juga
($\vect{AC} = 4\vect{AB}$).
\end{solution}

\begin{solution}{exo:g10:vectors:8}
$\vect{AB}\,(4, 2)$ dan $\vect{DC}\,(6-2,\ 1-(-1)) = (4, 2)$: kedua
\omterm{def:g10:vectors:vector}{vektornya} \emph{sama}, jadi $ABCD$ adalah jajargenjang, dan khususnya
$(AB) \parallel (CD)$. Sifat jajargenjang itulah yang lebih kuat: sifat itu
mengakibatkan kesejajarannya (kekolinearan \omterm{def:g10:vectors:vector}{vektornya}), tidak sebaliknya.
\end{solution}

\begin{solution}{exo:g10:vectors:9}
\emph{1.} $\vect{AB}\,(6, 3)$, sehingga
$\vect{AI} = \frac23\vect{AB} = (4, 2)$ dan $I(4, 2)$.

\emph{2.} $\vect{CB}\,(6-2,\ 3-5) = (4, -2)$, sehingga
$\vect{CJ} = \frac23\vect{CB} = \left(\frac83, -\frac43\right)$ dan
$J\left(2 + \frac83,\ 5 - \frac43\right) = \left(\frac{14}{3},
\frac{11}{3}\right)$.

\emph{3.} $\vect{IJ}\,\left(\frac{14}{3} - 4,\ \frac{11}{3} - 2\right) =
\left(\frac23, \frac53\right)$ dan $\vect{AC}\,(2, 5)$. Kriteria
kekolinearannya:
\[
\frac23 \times 5 - \frac53 \times 2 = \frac{10}{3} - \frac{10}{3} = 0 :
\]
kedua \omterm{def:g10:vectors:vector}{vektornya} \omterm{def:g10:vectors:collinear}{kolinear}, jadi $(IJ)$ sejajar dengan $(AC)$. (Baik $I$
maupun $J$ terletak sepertiga jalan dari sisinya menuju $B$: segitiga
kecil $IBJ$ adalah perkecilan $ABC$.)
\end{solution}

\begin{solution}{exo:g10:vectors:10}
\omterm{prop:g10:coordgeom:midpoint}{Titik tengahnya} adalah
\[
I\left(\tfrac{x_A + x_B}{2}, \tfrac{y_A + y_B}{2}\right), \quad
J\left(\tfrac{x_B + x_C}{2}, \tfrac{y_B + y_C}{2}\right), \quad
K\left(\tfrac{x_C + x_D}{2}, \tfrac{y_C + y_D}{2}\right), \quad
L\left(\tfrac{x_D + x_A}{2}, \tfrac{y_D + y_A}{2}\right).
\]
\omterm{def:g10:coordgeom:system}{Koordinat} pertama kedua \omterm{def:g10:vectors:vector}{vektornya}:
\[
x_J - x_I = \frac{x_B + x_C}{2} - \frac{x_A + x_B}{2}
= \frac{x_C - x_A}{2},
\qquad
x_K - x_L = \frac{x_C + x_D}{2} - \frac{x_D + x_A}{2}
= \frac{x_C - x_A}{2},
\]
dan perhitungan yang sama berlaku untuk \omterm{def:g10:coordgeom:system}{koordinat} keduanya. Jadi
$\vect{IJ} = \vect{LK}$ (keduanya sama dengan $\frac12\vect{AC}$), dan
$IJKL$ adalah jajargenjang, bagaimanapun bentuk segi empat $ABCD$.
\end{solution}

\begin{solution}{pb:g10:vectors:1}
\textbf{1.} $\vect{AB} + \vect{BC} + \vect{CD} = \vect{AD}$;
\quad $\vect{MA} - \vect{MB} = \vect{MA} + \vect{BM} =
\vect{BA}$; \quad $\vect{AB} + \vect{BC} + \vect{CD} +
\vect{DA} = \vect{AA} = \vec 0$ (setelah \omterm{def:g10:vectors:sum}{jumlah} ketiga diurutkan ulang).

\textbf{2.} $\vect{OB} - \vect{OA} = \vect{AO} + \vect{OB} =
\vect{AB}$ menurut Chasles.

\textbf{3.} $I$ adalah \omterm{prop:g10:coordgeom:midpoint}{titik tengah} $[AB]$ tepat ketika
$\vect{IA} = -\vect{IB}$, yaitu
$\vect{IA} + \vect{IB} = \vec 0$. Lalu, untuk sebarang $O$:
$\vec 0 = \vect{IO} + \vect{OA} + \vect{IO} + \vect{OB}$, sehingga
$2\,\vect{OI} = \vect{OA} + \vect{OB}$ --- dan sebaliknya.

\textbf{4.} $\vect{AB} = (4, 2)$, dan $D$ memenuhi
$\vect{DC} = \vect{AB}$: $D = C - (4, 2) = (2, -2)$ --- yaitu
$D$ yang sama dengan \cref{pb:g10:coordgeom:1}, satu baris lebih cepat.

\textbf{5.} Translasi oleh \omterm{def:g10:vectors:vector}{vektor} $\vec u$ --- persis
transformasi yang dihasilkan oleh dua cermin sejajar atau dua
setengah putaran. \omterm{def:g10:functions:function}{Peta} $(1, 2)$: $(4, 1)$.

\textbf{6.} Dengan \omterm{def:g10:algebra:expand}{menjabarkan} setiap $\vect{GX}$ sebagai
$\vect{GO} + \vect{OX}$:
$3\,\vect{GO} + \vect{OA} + \vect{OB} + \vect{OC} = \vec 0$,
sehingga $\vect{OG} = \frac13(\vect{OA} + \vect{OB} + \vect{OC})$:
hanya ada satu titik semacam itu, dan \omterm{def:g10:coordgeom:system}{koordinatnya} adalah rata-rata
\omterm{def:g10:coordgeom:system}{koordinat} ketiga titik sudutnya.

\textbf{7.} $\vect{AB} + \vect{AC} =
(\vect{AG} + \vect{GB}) + (\vect{AG} + \vect{GC})
= 2\vect{AG} + (\vect{GB} + \vect{GC})
= 2\vect{AG} - \vect{GA} = 3\vect{AG}$, sehingga
$\vect{AG} = \frac13(\vect{AB} + \vect{AC})$. Rumus \omterm{prop:g10:coordgeom:midpoint}{titik
tengah} (pertanyaan 3, dilihat dari $A$) memberi
$\vect{AK} = \frac12(\vect{AB} + \vect{AC})$. Setelah dibandingkan:
$\vect{AG} = \frac23\vect{AK}$ --- jadi $G$ terletak pada garis berat
$[AK]$, pada dua pertiga jalan dari $A$.

\textbf{8.} Syarat pendefinisinya memperlakukan $A$, $B$, $C$
secara sama, sehingga perhitungan yang sama dari $B$ atau dari $C$
menempatkan $G$ yang \emph{sama} pada dua pertiga garis berat itu:
perpotongan di satu titik tidak menuntut ongkos tambahan. Bukti
jajargenjangnya cerdik tetapi khusus; \omterm{def:g10:coordgeom:system}{koordinat} menghitung tetapi
mengaburkan; \omterm{def:g10:vectors:vector}{vektor} membuat kesetangkupannya terlihat dan buktinya dua baris.

\textbf{9.} $G\left(\frac{1+5+3}{3}, \frac{1+2+6}{3}\right) =
(3, 3)$; $K(4, 4)$; $\vect{AG} = (2, 2)$ dan
$\vect{AK} = (3, 3)$: determinannya
$2 \times 3 - 2 \times 3 = 0$, jadi \omterm{def:g10:vectors:collinear}{kolinear}, dan memang
$\vect{AG} = \frac23\vect{AK}$.

\textbf{10.} $G'\left(\frac{2 + 5 + 3}{4}, \frac{2 + 2 +
6}{4}\right) = (2.5,\ 2.5)$: tertarik dari $G(3,3)$ menuju
titik sudut $A(1,1)$ yang massanya digandakan, sebagaimana seharusnya
sebuah titik seimbang.

\textbf{11.} $3 \times 4 - (-2)(-6) = 12 - 12 = 0$: \omterm{def:g10:vectors:collinear}{kolinear}
($\vec v = -2\vec u$). $2 \times 9 - 5 \times 4 = -2 \neq 0$:
tidak \omterm{def:g10:vectors:collinear}{kolinear}.

\textbf{12.} $\vect{AB} = (2, 3)$, $\vect{AC} = (6, 9)$:
determinannya $2 \times 9 - 3 \times 6 = 0$: segaris.

\textbf{13.} $\vect{OI} = \frac12(\vect{OA} + \vect{OB})$ dan
$\vect{OK} = \frac12(\vect{OC} + \vect{OD})$, sehingga \omterm{prop:g10:coordgeom:midpoint}{titik tengah}
$[IK]$ mempunyai \omterm{def:g10:vectors:vector}{vektor} posisi
$\frac12(\vect{OI} + \vect{OK}) = \frac14(\vect{OA} +
\vect{OB} + \vect{OC} + \vect{OD})$ --- dan perhitungan
untuk $[JL]$ memberi \omterm{def:g10:vectors:vector}{vektor} yang persis sama. Kedua bimedian itu (dan,
menurut Varignon, kedua diagonal jajargenjang $IJKL$) berbagi
pusat ini: yaitu titik berat segi empatnya.

\textbf{14.} $\vect{MN} = \vect{AN} - \vect{AM} =
\frac13\vect{AC} - \frac13\vect{AB} = \frac13\vect{BC}$:
jadi $(MN) \parallel (BC)$ dan $MN = \frac13 BC$.

\textbf{15.} Dengan $\vect{AM} = k\,\vect{AB}$ dan
$\vect{AN} = k\,\vect{AC}$: $\vect{MN} = k\,\vect{BC}$ --- satu
baris yang memuat teorema \omterm{prop:g10:coordgeom:midpoint}{titik tengah} ($k = \frac12$) dan
Thales (untuk sebarang $k$): perbandingannya lolos utuh melalui
pengurangan itu.

\textbf{16.} Resultannya $(3, 4)$, besarnya
$\sqrt{9 + 16} = 5$ km/jam: perahu itu sesungguhnya bergerak
ke timur laut condong ke utara, sepanjang sisi miring segitiga
$3$--$4$--$5$.

\textbf{17.} $\vec u + \vec v + \vec w = (0, 0)$: cincin itu
diam --- seimbang. Sepanjang segi banyak tertutup
$A_1 A_2 \dots A_n A_1$, Chasles melipat \omterm{def:g10:vectors:sum}{jumlahnya} menjadi
$\vect{A_1 A_1} = \vec 0$: berkeliling membawamu kembali ke tempat semula.

\textbf{18.} $\vec w = -(\vec u + \vec v) = -(3, 2) =
(-3, -2)$.

\textbf{19.} Resultannya $(1.5,\ 2)$, besarnya
$\sqrt{2.25 + 4} = 2.5$ m/s. Penyeberangannya: $50$ m pada $2$ m/s
melintang: $25$ s; hanyutnya: $1.5 \times 25 = 37.5$ m ke hilir.

\textbf{20.} Yang dibuktikan ulang akhir pekan ini: teorema dua
pertiga titik berat (pertanyaan 7), pusat bimedian dan Varignon
(pertanyaan 13), serta Thales dan \omterm{prop:g10:coordgeom:midpoint}{titik tengah} (pertanyaan 14--15). Yang
masih belum dapat dilakukan anak panah: mengukur \emph{sudut} (dan
karenanya panjang pada arah sebarang) --- hasil kali skalar tahun
berikutnya justru merupakan alat yang hilang itu.

\end{solution}
